\chapter*{General Conclusion}
\markboth{General Conclusion}{}
\addcontentsline{toc}{chapter}{Conclusion Générale}

This project addressed the challenge of knowledge management at Jade
Advisory by designing and implementing Ask-Jade, an intelligent
question-answering system based on a Retrieval-Augmented Generation
architecture. By transforming the organization's dispersed document
library into a centralized and searchable knowledge base, the system
empowers consultants to access relevant project information faster and
more precisely, streamlining processes that were once time-consuming
and manually intensive.

Extensive research was conducted to identify and evaluate large language
models and embedding models suited to the project's requirements.
Following a comparative analysis, BGE-M3 was selected as the embedding
model for its multilingual capabilities and hybrid retrieval support,
while a configurable LLM layer allows the administrator to switch
generation models directly from the interface without any redeployment.
LlamaIndex was chosen as the orchestration framework for its native
optimization for document indexing and its compatibility with Docling
and pgvector, which form the core of the ingestion pipeline.

The evaluation results confirm the relevance of the proposed approach.
The final vector-only pipeline achieves strong performance across
retrieval and generation metrics, with a hallucination rate below the
target threshold and significant improvements in faithfulness and
helpfulness compared to the initial graph-based architecture. The
decision to remove the Neo4j knowledge graph was justified by both
quantitative results and practical considerations of cost, latency,
and extraction reliability on multilingual content.

Beyond its technical contribution, Ask-Jade provides Jade Advisory
with a secure and scalable solution that preserves data privacy,
supports multilingual content in English, French, and Arabic, and
gives users full traceability over generated answers through source
citations linked to the original document chunks.

As future work, several directions can be explored, including extending
the evaluation benchmark, integrating an Excel lookup tool to allow
consultants to query and extract structured data directly from
spreadsheet-based project records, and developing a document management
interface that centralizes the organization, versioning, and lifecycle
tracking of all ingested files within the platform.